\documentclass[10pt,letterpaper]{article}
\usepackage{cornell-notes}

\title{Clause 07.06.06: Additive Operators}
\author{}
\date{}
\setDocTitle{Clause 07.06.06: Additive Operators}
\setDocSubtitle{C++ 2024}
\setDocOwner{C++ 2024 Cornell Notes}
\setDocDate{\today}
\setCornellCollection{C++ 2024 Cornell Notes}
\setCornellUnitType{Clause}
\setCornellUnitNumber{07.06.06}
\setCornellUnitTitle{Additive Operators}
\hypersetup{pdftitle={Clause 07.06.06: Additive Operators},pdfsubject={Cornell notes},pdfauthor={}}

\begin{document}
\maketitle
\begin{CornellOverviewBox}{Overview}
\textbf{Classification:} Core language - Normative\\[2pt]
\textbf{Learning objective:} Understand the rules, terminology, and semantic consequences associated with Additive operators.
\end{CornellOverviewBox}\begin{CornellSummaryBox}{Summary}
{\sffamily\bfseries\color{Accent} Summary}\par\vspace{3pt}Section 7.6.6 focuses on Additive operators. Use the listed grammar productions together with the semantic constraints and disambiguation rules referenced by the section. When applying it, preserve the distinction between mandatory requirements, recommendations, permissions, and behavior deliberately left undefined, unspecified, or implementation-defined.
\end{CornellSummaryBox}

\end{document}
